% Template for producing ESWA-format journal articles using LaTeX
% Written by Miha Ravber
% Programming methodologies laboratory
% Faculty of Electrical Engineering and Computer Science
% University of Maribor
% Koroška cesta 46, 2000 Maribor
% E-mail: miha.ravber@um.si
% WWW: https://lpm.feri.um.si/en/members/ravber/
% Created: November 20, 2020 by Miha Ravber
% Modified: November 9, 2023 by Miha Ravber
% Modified: February 21, 2024 by Miha Ravber
% Use at your own risk :)
% Please submit your issues on the github page: https://github.com/Ravby/eswa-template


\documentclass[review]{elsarticle}
\graphicspath{ {./figures/} }
\usepackage{hyperref}
\usepackage{float}
\usepackage{verbatim} %comments
\usepackage{apalike}
\restylefloat{figure}
\floatstyle{plaintop} %table caption at top
\restylefloat{table}
\usepackage[utf8]{inputenc}
\usepackage[T1]{fontenc}
\usepackage{amsmath,amssymb,amsthm}
\usepackage{bm}
\usepackage{graphicx}
\usepackage[margin=1in]{geometry}
\usepackage{enumitem}
\usepackage{booktabs}
\usepackage{multirow}
\usepackage{tikz}
\usetikzlibrary{calc}
\usetikzlibrary{arrows.meta, positioning, fit, backgrounds}
\usetikzlibrary{positioning,shapes.geometric,fit,backgrounds,calc,arrows.meta}
\usepackage[mode=tex,subpreambles=true]{standalone}
%% ------------------------------------------------------------------
%% REVISION MARKUP: methodology changes driven by the implementation.
%% Red text marks what changed relative to the previous draft; a red
%% footnote beside each substantive change explains what changed and
%% why (draft copy for co-author review). The same explanations sit in
%% this source as comments beginning with "%% CHANGED". Flip
%% \showrevfalse below to print a clean version: all red reverts to
%% normal ink and every change-note footnote disappears.
%% ------------------------------------------------------------------
\usepackage{xcolor}
\newif\ifshowrev
\showrevtrue
\ifshowrev
  \newcommand{\rev}[1]{\textcolor{red}{#1}}
  \newcommand{\revnote}[2]{\textcolor{red}{#1}\footnote{\textcolor{red}{[change note]} #2}}
\else
  \newcommand{\rev}[1]{#1}
  \newcommand{\revnote}[2]{#1}
\fi

\tikzset{
  datanode/.style={rectangle, rounded corners=2pt, draw=black!60, fill=black!8,
    minimum width=3.6cm, minimum height=1.05cm, align=center, font=\small},
  modelnode/.style={rectangle, rounded corners=2pt, draw=violet!70!black, fill=violet!12,
    minimum width=3.4cm, minimum height=1.05cm, align=center, font=\small},
  econnode/.style={rectangle, rounded corners=2pt, draw=teal!70!black, fill=teal!12,
    minimum width=3.4cm, minimum height=1.05cm, align=center, font=\small},
  outnode/.style={rectangle, rounded corners=2pt, draw=teal!70!black, fill=teal!22,
    minimum width=3.6cm, minimum height=1.05cm, align=center, font=\small},
  layerbox/.style={rectangle, rounded corners=4pt, draw=black!40, dashed, inner sep=8pt},
  flow/.style={-{Stealth[length=2.5mm]}, thick, black!70}
}

\newtheorem{problem}{Problem}
\newcommand{\R}{\mathbb{R}}
\newcommand{\E}{\mathbb{E}}
\newcommand{\Prob}{\mathbb{P}}
\journal{Expert Systems with Applications}

%% For ESWA journal you need to use APA style
\bibliographystyle{model5-names}\biboptions{authoryear}

\begin{document}
\begin{frontmatter}

%The title page must contain the title of the paper and the full name/full affiliation with country/e-mail address for each author and co-author of the manuscript. Please make sure you have included all elements listed below with your manuscript submission.

\begin{titlepage}
\begin{center}
\vspace*{1cm}

\textbf{ \large Expert Systems with Applications \LaTeX\ template}

\vspace{1.5cm}

% Author names and affiliations
First Author$^{a,b}$ (first.author@mail.com), Second Author$^a$ (second.author@mail.com), Last Author$^c$ (last.author@mail.com) \\

\hspace{10pt}

\begin{flushleft}
\small
$^a$ Full address of first author, including the country name \\
$^b$ Full address of second author, including the country name \\
$^c$ Full address of last author, including the country name

\begin{comment}
Clearly indicate who will handle correspondence at all stages of refereeing and publication, also post-publication. Ensure that phone numbers (with country and area code) are provided in addition to the e-mail address and the complete postal address. Contact details must be kept up to date by the corresponding author.
\end{comment}

\vspace{1cm}
\textbf{Corresponding author at: Full address of the corresponding author, including the country name.} \\
Last Author \\
Full address of the corresponding author, including the country name \\
Tel: (555) 555-1234 \\
Email: last.author@mail.com

\end{flushleft}
\end{center}
\end{titlepage}

\title{Predictive Modeling of Cyber Threats and Systemic Economic Risk in the United States: A Bayesian Hierarchical State-Space Approach}

\author[label1,label2]{First Author}
\ead{first.author@mail.com}

\author[label1]{Second Author}
\ead{second.author@mail.com}

\author[label3]{Last Author \corref{cor1}}
\ead{last.author@mail.com}

\cortext[cor1]{Corresponding author.}
\address[label1]{Full address of first author, including the country name}
\address[label2]{Full address of second author, including the country name}
\address[label3]{Full address of last author, including the country name}

\begin{abstract}
Cyber threats have evolved from isolated technical incidents into a source of systemic economic risk, capable of propagating through interdependent digital and economic infrastructures. Yet the literatures that study these two phenomena, technical threat prediction on the one hand and economic loss and systemic-risk modeling on the other, remain largely disconnected: forecasting stops at counts and scores, actuarial models begin from losses already realized, and macroeconomic models begin from scenarios someone invented. This paper proposes a unified Bayesian hierarchical state-space framework that closes that gap. Sector-level cyber threat intensities in the United States are modeled as latent, partially observed stochastic processes, and heterogeneous public data streams, namely vulnerability disclosures, exploitation intelligence and mandatory incident disclosures, are treated not as measurements of those intensities but as noisy observation channels, each with its own reporting bias. Cross-threat dependence is carried by a small set of shared dynamic factors and abrupt structural breaks by a latent regime, so that co-movement and regime change are inferred rather than assumed away. Probabilistic threat forecasts are then propagated through a software-to-sector exposure map and inter-sectoral economic linkages, yielding full predictive distributions of aggregate economic exposure in which uncertainty from every layer, from the latent threat state through the damage functions to the input--output propagation, is carried forward rather than discarded. Because the outputs are distributions rather than point estimates, they support value-at-risk, expected-shortfall and threshold-exceedance statements, and they are evaluated out of sample with proper scoring rules and calibration diagnostics, addressing methodological gaps that persist across the vulnerability-prediction, actuarial and systemic-risk literatures.
\end{abstract}

\begin{keyword}
ESWA \sep \LaTeX \sep Elsevier \sep template
\end{keyword}

\end{frontmatter}

\section{Introduction}
\label{introduction}

Modern economies run on digital infrastructure that is under sustained attack, and the United States (the largest, most digitally intermediated economy) concentrates a disproportionate share of global cyber exposure: scenario analysis of a single systemic event on a major payments system attributes the largest five-year national loss, on the order of one trillion USD, to the US alone \citep{lloyds2023}. Individual incidents already show the path from technical compromise to economic disruption: the 2017 NotPetya campaign halted global logistics, with estimated damages of \$3--57 billion depending on modeling assumptions \citep{welburn2022}; the 2021 Colonial Pipeline ransomware attack interrupted East Coast fuel distribution; and the 2024 Change Healthcare breach disrupted payments across much of US healthcare. Two features make the problem hard. Cyber risk is \emph{non-stationary and adversarial} (vulnerabilities, exploitation technology, and attacker incentives all evolve, shifting the data-generating process over time) and it is \emph{systemic}: shared software monocultures, cloud concentration, supply chains, and financial networks couple firms so that a local technical shock cascades sector- and economy-wide \citep{esrb2020,eisenbach2022}. Anticipating it therefore demands not the next vulnerability or insurance price but a forward-looking, probabilistic mapping from the technical threat landscape to aggregate economic exposure (which no peer-reviewed work provides).

Answering that demand means answering four questions at once, and each has been studied largely on its own. Which vulnerability classes and attack vectors will dominate, and at what intensity: the arrival process of disclosures, new attack topics, and weaponization. What those threats cost once realized: the frequency and severity of losses and their dependence across firms, sectors, and threat categories, where correlated and self-exciting arrivals break the independence assumptions behind conventional loss models and diversification. How a shock in one firm or sector reaches the rest of the economy, through inter-sectoral input--output linkages, payment and settlement networks, shared service providers, and confidence effects. And how any of it can be estimated from censored, delayed, biased data, given that vulnerability counts confound true threat activity with discovery effort and administrative backlogs, incidents are chronically under-reported, and the process exhibits abrupt regime change, so that a credible model must quantify rather than suppress the resulting uncertainty.

The forecasting literature answers the first question and stops there. \citet{kia2024} propose \textsc{CyRiPred}, mapping CVE free-text descriptions to Wikipedia-derived cybersecurity topics and aggregating topic occurrence and CVSS severity into monthly risk-score series forecast by supervised regressors (random forests best); \citet{subroto2019} combine CVE records with Twitter discussion to classify emerging risks; \citet{leverett2020} forecast aggregate CVE volume; and EPSS estimates each published CVE's probability of in-the-wild exploitation within 30 days via gradient boosting on exploitation telemetry \citep{jacobs2021,jacobs2023}. These are genuinely predictive and data-driven but terminate at the technical layer (counts, scores, probabilities); none translates threat intensity into economic quantities, virtually all issue \emph{point} forecasts without calibrated predictive distributions, and most treat each series univariately (a limitation \citet{kia2024} concede, leaving cross-topic correlation to future multivariate work).

The actuarial literature answers the second. \citet{eling2019} classify operational cyber losses and fit frequency--severity distributions; \citet{xu2018} model hacking-breach inter-arrival times and sizes with copula dependence; \citet{eling2018} and \citet{peng2018} employ vine copulas for cross-sectional dependence among breach losses; \citet{bessyroland2021} show multivariate Hawkes processes capture the self- and cross-excitation of breaches by type and target; and \citet{awiszus2023} unify idiosyncratic, systematic, and systemic risk for pricing. This stream has the probabilistic sophistication the fourth question demands, but it is calibrated \emph{backward} on realized incidents and decoupled from any forward-looking technical signal; \citet{carannante2023} further note that dependence among risk classes remains underexplored and that ignoring it distorts capital requirements.

The macroeconomic literature answers the third. \citet{welburn2022} extend sector-level input--output analysis to cyber via cascading, common-cause, and independent failures; \citet{eling2023} give a standardized framework for extreme-scenario macroeconomic cost; \citet{eisenbach2022} pre-mortem an attack on the US wholesale payments network, estimating an attack on any top-five bank impairs on average 31\% of the network, with forgone payments reaching multiples of daily GDP under liquidity hoarding; the ESRB formalizes the shock--amplification--systemic-event chain \citep{esrb2020}; and industry studies reach trillion-dollar magnitudes \citep{lloyds2023}. But the shocks driving these models are \emph{hypothetical}, posited by experts or drawn from stylized scenarios rather than forecast, so estimates carry order-of-magnitude uncertainty (the \$3--57 billion NotPetya range is emblematic \citep{welburn2022}) and, being scenarios rather than forecasts, are never validated out of sample.

The fourth question is answered only sporadically and locally: state-space models for weekly attack counts within single organizations \citep{condon2008,okutan2018}, Bayesian networks for enterprise risk assessment, and Bayesian updating of threat-category probabilities from intelligence feeds. None provides hierarchical, fully probabilistic inference across the threat--loss--economy pipeline, treats reporting bias as an explicit measurement model, or accommodates regime change at the national scale.

The gaps follow from this separation. Technical models stop at scores, actuarial models start from realized losses, and macroeconomic models start from invented scenarios, so nothing carries a forecast of the threat landscape through to an economic quantity. The predictive stream reports point forecasts and error metrics (MSE, MAE) rather than calibrated predictive distributions, so tail risk (what matters for systemic analysis, insurance capital, and policy) cannot be assessed. Category series are forecast independently despite evident co-movement and contagion, and the dependence machinery developed for realized losses has not been applied to forward-looking series. Models assume a stable data-generating process despite abrupt structural breaks (ransomware industrialization, Log4Shell, AI-assisted exploitation). CVE counts are treated as direct measurements of threat activity rather than noisy, effort-dependent observations of a latent process, and incident data are used without correcting under-reporting and disclosure lags. Scenario-based systemic outputs are unfalsifiable, and the field lacks probabilistic forecast evaluation with proper scoring rules.

Our contributions are as follows. An end-to-end probabilistic pipeline linking vulnerability-derived threat dynamics, through a software-to-sector exposure mapping, to inter-sectoral economic propagation for the US, replacing hypothetical scenario shocks with data-driven forecasts. A hierarchical Bayesian state-space formulation yielding exact finite-sample posterior predictive distributions over threat intensities and losses (enabling predictive value-at-risk and tail-expectation statements) and pooling strength across thousands of sparse topics via hierarchical shrinkage. Multivariate cross-threat co-movement and contagion via a dynamic latent factor structure with a self-exciting extension, directly answering the multivariate extension \citet{kia2024} call for. Latent Markov-switching states that detect and adapt to structural breaks, with posterior regime probabilities as an interpretable early-warning signal. An explicit measurement model treating CVE publications, exploitation signals (EPSS/KEV), and mandatory disclosures (SEC Form 8-K, Item 1.05) as distinct, bias-prone observation channels of a common latent threat state, exploiting a disclosure regime in force since December 2023 that predates nearly all existing models. All outputs issued as predictive distributions and evaluated out of sample against realized vulnerability, exploitation, and loss data using proper scoring rules (CRPS, log score) and calibration diagnostics, establishing a validation standard for systemic cyber-economic modeling.


\section{Problem Formulation}\label{sec:formulation}

Time is discrete, $t = 1, \dots, T$ (months). Let $\mathcal{K} = \{1, \dots, K\}$ index cyber threat categories (e.g., topics extracted from vulnerability descriptions, or curated classes such as ransomware, supply-chain compromise, injection attacks) and $\mathcal{S} = \{1, \dots, S\}$ index US economic sectors (e.g., BEA/NAICS summary industries). The model's central object is the \emph{latent threat intensity} $\bm{\lambda}_t = (\lambda_{1,t}, \dots, \lambda_{K,t})^\top \in \R_{+}^{K}$, where $\lambda_{k,t}$ is the true, unobserved level of malicious activity and exploitable exposure for category $k$ in period $t$. It is not directly observable; we observe only noisy, biased channels: \textbf{vulnerability disclosures} $N_{k,t} \in \mathbb{Z}_{\geq 0}$, the count of CVE records published in $t$ assigned to $k$ (via a topic-extraction map in the spirit of \textsc{WikiTopic} \citep{kia2024}), with severity marks $B_{k,t} \in [0,10]$ (mean CVSS base score); \textbf{exploitation signals} $E_{k,t} \in [0,1]$, an aggregate exploitation-likelihood signal for $k$ (e.g., mean or upper-quantile EPSS score, CISA Known Exploited Vulnerabilities inclusion counts); and \textbf{incident disclosures} $D_{s,t} \in \mathbb{Z}_{\geq 0}$, the count of material cyber incidents disclosed by sector-$s$ firms in $t$ (SEC 8-K Item 1.05 filings and curated repositories), with loss marks $L_{s,t} \geq 0$ where available. Economic exposure to the threat landscape is a (possibly time-varying) map $\bm{M}_t \in [0,1]^{S \times K}$, $M_{s,k,t} \approx$ the share of sector $s$'s digital footprint exposed to category $k$, built from affected-product metadata (NVD CPE fields) matched to sectoral technology-adoption data. Inter-sectoral linkages follow the US input--output structure: $\bm{A} \in \R_{+}^{S \times S}$ is the technical-coefficients matrix from the BEA Use tables and $\bm{\Lambda}_L = (\bm{I} - \bm{A})^{-1}$ its Leontief inverse.

Following the risk-as-likelihood-times-impact convention \citep{kia2024}, but lifted to the economic layer, the sectoral cyber shock is
\begin{equation}\label{eq:shock}
    g_{s,t} \;=\; \phi_s\!\Big( \textstyle\sum_{k=1}^{K} M_{s,k,t}\, \lambda_{k,t}\, \sigma_{k,t} \Big),
\end{equation}
with $\sigma_{k,t}$ the latent severity of $k$ at $t$ and $\phi_s(\cdot)$ a sector-specific damage function mapping technical shock load to the fraction of sectoral output disrupted, calibrated on documented historical events. Direct sectoral losses $\bm{d}_t = (g_{1,t} x_{1,t}, \dots, g_{S,t} x_{S,t})^\top$, with $x_{s,t}$ sectoral gross output, propagate to total (direct plus indirect) losses via
\begin{equation}\label{eq:leontief}
    \bm{\ell}_t \;=\; \bm{\Lambda}_L\, \bm{d}_t,
    \qquad
    \ell_t^{\mathrm{agg}} \;=\; \bm{1}^\top \bm{\ell}_t .
\end{equation}

The following three problems are thus posed. \textbf{(1) Probabilistic threat forecasting:} given the observation history $\mathcal{F}_t = \sigma\big(\{N_{k,\tau}, B_{k,\tau}, E_{k,\tau}, D_{s,\tau}, L_{s,\tau}\}_{\tau \le t}\big)$, characterize the posterior predictive $p\big(\bm{\lambda}_{t+1:t+h} \mid \mathcal{F}_t\big)$, $h \geq 1$, including cross-category dependence and regime uncertainty, rather than point forecasts $\hat{\bm{\lambda}}_{t+h}$. \textbf{(2) Systemic economic exposure:} characterize the induced posterior predictive of aggregate loss $p\big(\ell^{\mathrm{agg}}_{t+1:t+h} \mid \mathcal{F}_t\big)$ and its tail functionals (predictive value-at-risk $\mathrm{VaR}_{\alpha}$ and expected shortfall $\mathrm{ES}_{\alpha}$ of $\ell^{\mathrm{agg}}_{t+h}$, and the systemic-event probability $\Prob(\ell^{\mathrm{agg}}_{t+h} > c \mid \mathcal{F}_t)$ for policy-relevant thresholds $c$). \textbf{(3) Inference under biased observation:} perform (1)--(2) while acknowledging that $N_{k,t}$ confounds true activity with discovery effort and administrative throughput, $D_{s,t}$ under-counts incidents with sector- and time-varying reporting propensity, and the latent dynamics undergo structural breaks; formally, the channels are conditionally independent given the latent state, each with its own bias process, and the latent state follows regime-switching dynamics.


\section{Proposed Solution}\label{sec:solution}
The proposed framework is a single Bayesian hierarchical state-space model in which public data channels (NVD CVE+CPE vulnerability data $N_{k,t}, B_{k,t}$; EPSS/CISA KEV exploitation data $E_{k,t}$; SEC 8-K incident filings $D_{s,t}, L_{s,t}$) are treated as biased observations of latent, regime-switching threat dynamics, and posterior predictive draws are propagated through economic translation (Figure~\ref{fig:pipeline}). This section states the framework as a procedure running from raw public data to the four delivered outputs.

%% ------------------------------------------------------------------
%% Pipeline figure, drawn inline. Nothing to upload, no external file.
%% Uses only libraries main.tex already loads, plus the paper's own \E
%% and \Prob. Styles scoped so they cannot leak into datanode etc.
%% ------------------------------------------------------------------
\begingroup
%% elsarticle[review] is double spaced; \baselinestretch also applies inside
%% tikz nodes, which would make every multi-line node taller than designed.
\renewcommand{\baselinestretch}{1}\selectfont
\providecommand{\ttl}[1]{\textcolor{blue!85!black}{\textbf{#1}}}
\begin{figure}[p]
\centering
\tikzset{
  base/.style   = {rounded corners=4pt, draw=black!45, line width=0.7pt,
                   align=center, inner sep=4pt, font=\fontsize{7}{8}\selectfont},
  proc/.style   = {base, fill=white},
  data/.style   = {base, fill=cyan!25},
  outp/.style   = {base, fill=pink!55},
  optn/.style   = {base, fill=white, dashed},
  dec/.style    = {diamond, aspect=2.1, draw=black!45, fill=white,
                   align=center, inner sep=1pt, font=\fontsize{7}{8}\selectfont},
  cbox/.style   = {rounded corners=9pt, draw=black!45, fill=black!5,
                   line width=0.9pt, inner sep=8pt},
  lab/.style    = {font=\fontsize{9}{11}\selectfont\bfseries, anchor=south west, xshift=5pt, yshift=3pt},
  fl/.style     = {-{Stealth[length=5pt,width=4pt]}, draw=black!65, line width=0.7pt},
  fld/.style    = {-{Stealth[length=5pt,width=4pt]}, draw=black!50, line width=0.7pt, dashed},
  elab/.style   = {font=\fontsize{7}{8}\selectfont, inner sep=1.5pt, fill=white},
  W/.style      = {text width=5.85cm},
  H/.style      = {text width=2.6cm},
}
\resizebox{\textwidth}{!}{%
\begin{tikzpicture}[node distance=0.26cm and 0.26cm]

%% =====================================================================
%% ROW 1 : data -> measurement -> latent dynamics -> inference
%% =====================================================================

% ---------- Group A : observed inputs --------------------------------
\node[data, text width=3.8cm] (cve) at (0,0)
  {\ttl{NVD: CVE $+$ CPE}\\[1pt] monthly counts $N_{k,t}$, severity marks $B_{k,t}$, CPE product metadata};
\node[data, below=of cve, text width=3.8cm] (epss)
  {\ttl{EPSS / CISA KEV}\\[1pt] exploitation signal $E_{k,t}$};
\node[data, below=of epss, text width=3.8cm] (sec)
  {\ttl{SEC 8-K filings}\\[1pt] incident counts $D_{s,t}$, loss marks $L_{s,t}$};
\node[proc, below=0.8cm of sec, text width=3.8cm] (assemble)
  {\ttl{Assemble inputs}\\[2pt]
   $\mathcal{K}=\{1,\dots,K\}$, $\mathcal{S}=\{1,\dots,S\}$\\[2pt]
   history $\mathcal{F}_t=\sigma\big(\{N,B,E,D,L\}_{\tau\le t}\big)$};

% ---------- Group B : Layer 2 measurement ----------------------------
\node[proc, W, anchor=north] (effort) at (7.179,0)
  {\ttl{Reporting effort $e_t$}\\[2pt]
   standardized smooth trend of aggregate $\log$ CVE volume (NVD throughput, backlog, secular growth)};
\node[proc, H, below=of effort.south west, anchor=north west] (mvul)
  {\ttl{Vulnerability}\\[2pt] $N_{k,t}\sim\mathrm{NegBin}\big(e_t\lambda_{k,t},\psi_k\big)$};
\node[proc, H, right=of mvul] (mexp)
  {\ttl{Exploitation}\\[2pt] $\operatorname{logit}E_{k,t}\sim\mathcal{N}\big(\alpha_k+\beta_k\log\lambda_{k,t},\varsigma_k^{2}\big)$};
\node[proc, H, below=of mvul] (minc)
  {\ttl{Incidents}\\[2pt] $D_{s,t}\sim\mathrm{Poisson}\big(\rho_s\Lambda_{s,t}+\pi_s\big)$,\\ $\Lambda_{s,t}=\sum_k M_{s,k,t}\lambda_{k,t}$};
\node[proc, H, right=of minc] (msev)
  {\ttl{Severity}\\[2pt] $B_{k,t}\sim\mathrm{LogNormal}$\\ $\big(\log\sigma_{k,t},\,\kappa_k^{2}\big)$};
\node[proc, W, below=of minc.south west, anchor=north west] (cindep)
  {\ttl{Masking and conditional independence}\\[2pt]
   empty category-months masked, never imputed; channels independent given $\{\bm{\lambda}_t,\bm{\sigma}_t\}$};
\begin{scope}[on background layer]
  \node[cbox, fit=(effort)(mvul)(mexp)(minc)(msev)(cindep),
        label={[lab]north west:Measurement model}] (layer2) {};
\end{scope}

% ---------- Group C : Layer 1 latent dynamics ------------------------
\node[proc, H, anchor=north west] (reg) at (12.594,0)
  {\ttl{Markov regimes}\\[2pt] $z_t\in\{1,\dots,R\}$, $\bm{\Pi}$; levels $\bm{\mu}_{z_t},\bm{\nu}_{z_t}$};
\node[proc, H, right=of reg] (fac)
  {\ttl{Dynamic factors}\\[2pt] $\bm{f}_t=\bm{\Phi}\bm{f}_{t-1}+\bm{w}_t$, $\bm{w}_t\sim\mathcal{N}(\bm{0},\bm{Q})$};
\node[proc, W, below=of reg.south west, anchor=north west] (int)
  {\ttl{Log-intensities}\\[2pt]
   $\bm{\eta}_t=\bm{\mu}_{z_t}+\bm{\Gamma}\bm{f}_t+\bm{u}_t$, \quad $\bm{\lambda}_t=\exp(\bm{\eta}_t)$};
\node[proc, W, below=of int] (sevl)
  {\ttl{Log-severities}\\[2pt]
   $\log\bm{\sigma}_t=\bm{\nu}_{z_t}+\bm{\Psi}\bm{h}_t$ \ (analogous low-dimensional latent process)};
\node[proc, W, below=of sevl] (idn)
  {\ttl{Identification and priors}\\[2pt]
   leading $r\times r$ block of $\bm{\Gamma},\bm{\Psi}$ lower triangular, positive diagonal;
   half-Normal scales bounded away from zero};
\node[proc, W, below=of idn] (ker)
  {\ttl{Learned temporal kernel}\\[2pt]
   $g_t=\sigma_g\sum_d w_d(h)u_{t-d}$, \ $w_d(h)=\mathrm{sigmoid}\big((h-|d|)/2\big)$;
   enters $\bm{\eta}_t$ as $a_kg_t$; posterior of $h$ gives the co-movement span};
\begin{scope}[on background layer]
  \node[cbox, fit=(reg)(fac)(int)(sevl)(idn)(ker),
        label={[lab]north west:Latent threat dynamics}] (layer1) {};
\end{scope}

% ---------- Group D : Layer 4 inference ------------------------------
\node[proc, W, anchor=north west] (marg) at (21.653,0)
  {\ttl{Regime marginalization}\\[2pt]
   HMM forward recursion integrates out $z_{1:T}$ inside the likelihood, jointly with all four channels};
\node[proc, W, below=0.75cm of marg] (nuts)
  {\ttl{NUTS sampling}\\[2pt]
   joint posterior $p\big(\{\bm{f}_t,\bm{\lambda}_t\},\bm{\Theta}\mid\mathcal{F}_T\big)$ over all continuous unknowns, no discrete updates
   (where $\bm{\Theta}$ collects the static parameters $\bm{\Phi},\bm{Q},\bm{\Gamma},\bm{\Psi},\bm{\Pi},\bm{\mu},\bm{\nu}$, the noise scales, and the channel parameters)};
\node[inner sep=0pt, minimum size=0pt] (nbfit) at ($(nuts.west)+(-0.55,0)$) {};
\node[dec, below=0.55cm of nuts] (conv)
  {\ttl{Chains converged?}\\[1pt] $\widehat{R}\approx1$, ESS adequate,\\ no divergences\\ (where $\widehat{R}$ is the\\ between-chain variance ratio)};
\node[proc, W, below=0.55cm of conv] (ffbs)
  {\ttl{FFBS regime paths}\\[2pt] posterior draws of $z_{1:T}$ recovered after sampling};
\begin{scope}[on background layer]
  \node[cbox, fit=(marg)(nuts)(conv)(ffbs)(nbfit),
        label={[lab]north west:Posterior inference}] (layer4) {};
\end{scope}

%% =====================================================================
%% ROW 2 : forecast -> economic translation -> outputs -> validation
%% =====================================================================
\coordinate (r2) at (-1.756,-8.95);

% ---------- Group E : Stage F forecast simulation --------------------
\node[proc, W, anchor=north west] (floop) at (r2)
  {\ttl{For every posterior draw $m$ and horizon $h=1,\dots,12$}};
\node[proc, H, below=of floop.south west, anchor=north west] (fz)
  {\ttl{Regime path}\\[2pt] $z^{(m)}_{T+h}\sim$\\ $\bm{\Pi}^{(m)}_{z^{(m)}_{T+h-1},\,\cdot}$};
\node[proc, H, right=of fz] (ff)
  {\ttl{Factor propagation}\\[2pt] $\bm{f}^{(m)}_{T+h}=\bm{\Phi}^{(m)}\bm{f}^{(m)}_{T+h-1}+\bm{w}^{(m)}_{T+h}$};
\node[proc, W, below=of fz.south west, anchor=north west] (fl)
  {\ttl{Latent state at the drawn regime}\\[2pt]
   $\bm{\lambda}^{(m)}_{T+h}=\exp(\bm{\eta}^{(m)}_{T+h})$, $\bm{\sigma}^{(m)}_{T+h}$;
   \ effort $e_t$ carried at its last observed value};
\node[dec, below=0.55cm of fl] (hdec) {\ttl{$h=12$?}};
\begin{scope}[on background layer]
  \node[cbox, fit=(floop)(fz)(ff)(fl)(hdec),
        label={[lab]north west:Forecast simulation}] (stageF) {};
\end{scope}

% ---------- Group F : Layer 3 economic translation -------------------
\node[data, W, anchor=north west] (struct) at ($(r2)+(8.503,0)$)
  {\ttl{Structural inputs}\\[2pt]
   exposure map $\bm{M}_t\in[0,1]^{S\times K}$ (CPE $\to$ sector); BEA $\bm{A}$,
   $\bm{\Lambda}_L=(\bm{I}-\bm{A})^{-1}$; gross output $x_{s,t}$};
\node[data, W, below=of struct] (calib)
  {\ttl{Calibration set}\\[2pt]
   NotPetya, Colonial Pipeline, MOVEit, Change Healthcare loss ranges
   $+$ aggregate annual US-loss anchor};
\node[proc, W, below=of calib] (dmg)
  {\ttl{Damage functions $\phi_s$}\\[2pt]
   zero-anchored sigmoid family, three global parameters, MCMC-calibrated: each loss range a $90\%$ lognormal interval};
\node[proc, H, below=of dmg.south west, anchor=north west] (shock)
  {\ttl{Sectoral shock}\\[2pt] $g_{s,t}=\phi_s\big(S_{s,t}\big)$,\\ $S_{s,t}=\sum_k M_{s,k,t}\lambda_{k,t}\sigma_{k,t}$};
\node[proc, H, right=of shock] (dir)
  {\ttl{Direct losses}\\[2pt] $\bm{d}_t=\big(g_{s,t}\,x_{s,t}\big)_{s=1}^{S}$};
\node[proc, W, below=of shock.south west, anchor=north west] (leon)
  {\ttl{Leontief propagation}\\[2pt]
   $\bm{\ell}_t=\bm{\Lambda}_L\bm{d}_t$, \quad $\ell^{\mathrm{agg}}_t=\bm{1}^{\!\top}\bm{\ell}_t$};
\node[proc, W, below=of leon] (unc)
  {\ttl{Joint uncertainty propagation}\\[2pt]
   $\phi_s$, $\bm{M}_t$ and $\bm{\lambda}_{t+h}$ propagated together, so $\mathrm{VaR}_\alpha$ and $\mathrm{ES}_\alpha$ carry every layer};
\begin{scope}[on background layer]
  \node[cbox, fit=(struct)(calib)(dmg)(shock)(dir)(leon)(unc),
        label={[lab]north west:Economic translation}] (layer3) {};
\end{scope}

% ---------- Group G : outputs ----------------------------------------
\node[outp, text width=3.4cm, anchor=north west] (o1) at ($(r2)+(16.703,-0.9)$)
  {\ttl{Threat forecasts}\\[2pt] predictive intensity by category with credible bands};
\node[outp, text width=3.4cm, below=0.55cm of o1] (o2)
  {\ttl{Sector exposure indices}\\[2pt] $\E[g_{s,t+h}\mid\mathcal{F}_t]$};
\node[outp, text width=3.4cm, below=0.55cm of o2] (o3)
  {\ttl{Aggregate loss distribution}\\[2pt]
   $\mathrm{VaR}_{0.95}$, $\mathrm{ES}_{0.95}$, $\Prob(\ell^{\mathrm{agg}}_{t+h}>c\mid\mathcal{F}_t)$};
\node[outp, text width=3.4cm, below=0.55cm of o3] (o4)
  {\ttl{Regime alerts}\\[2pt] $\Prob(z_t=r\mid\mathcal{F}_T)$ as early warning};

%% =====================================================================
%% ARROWS
%% =====================================================================
% inputs -> assemble -> Layer 2
\draw[fl] (cve.south)  -- (epss.north);
\draw[fl] (epss.south) -- (sec.north);
\draw[fl] (sec.south)  -- (assemble.north);
\draw[fl] (assemble.east) -- ++(0.55,0) |- (layer2.west);

% Layer 2 internals
\draw[fl] (effort.south) -- ++(0,-0.12) -| (mvul.north);
\draw[fl] (effort.south) -- ++(0,-0.12) -| (mexp.north);
\draw[fl] (mvul.south) -- (minc.north);
\draw[fl] (mexp.south) -- (msev.north);

% Layer 2 -> Layer 1 -> Layer 4
\draw[fl] (layer2.east) -- node[elab,above=1pt,pos=0.5,align=center]{biased views\\ of $\bm{\lambda}_t,\bm{\sigma}_t$} (layer1.west);
\draw[fl] (layer1.east) -- node[elab,above=1pt,pos=0.5,align=center]{state-space\\ likelihood} (layer4.west);

% Layer 1 internals
\draw[fl] (reg.south) -- ++(0,-0.12) -| ($(int.north)+(-1.2cm,0)$);
\draw[fl] (fac.south) -- ++(0,-0.12) -| ($(int.north)+(1.2cm,0)$);
\draw[fl] (int.south) -- (sevl.north);
\draw[fl] (sevl.south) -- (idn.north);
\draw[fl] (ker.north) -- (idn.south);

% Layer 4 internals
\draw[fl] (marg.south) -- (nuts.north);
\draw[fl] (nuts.south) -- (conv.north);
\draw[fl] (conv.south) -- node[elab,pos=0.45]{yes} (ffbs.north);
\coordinate (nb) at ($(nuts.west)+(-0.55,0)$);
\draw[fl] (conv.west) -| (nb) -- (nuts.west);
\node[elab, right=1pt] at ($(nb)!0.5!(nb |- conv.west)$) {no};

% row wrap : Layer 4 -> Stage F
\coordinate (wA) at ($(stageF.north)+(1.60,0)$);
\coordinate (wy) at ($(stageF.north)+(0,0.62)$);
\coordinate (wg) at ($(layer4.west)+(-1.08,0)$);
\coordinate (w1) at ($(layer4.south)+(0,-0.45)$);
\coordinate (w2) at (wg |- w1);
\coordinate (w3) at (wg |- wy);
\coordinate (w4) at (wA |- wy);
\draw[fl] (layer4.south) -- (w1) -- (w2) -- (w3) -- (w4) -- (wA);
\node[elab, above] at ($(w3)!0.5!(w4)$) {posterior draws $m=1,\dots,M$ carried forward into the forecast recursion};

% Stage F internals
\draw[fl] (floop.south) -- ++(0,-0.12) -| (fz.north);
\draw[fl] (floop.south) -- ++(0,-0.12) -| (ff.north);
\draw[fl] (fz.south) -- ++(0,-0.12) -| ($(fl.north)+(-1.2cm,0)$);
\draw[fl] (ff.south) -- ++(0,-0.12) -| ($(fl.north)+(1.2cm,0)$);
\draw[fl] (fl.south) -- (hdec.north);
\draw[fl] (hdec.west) -| ($(floop.west)+(-0.16,0)$) -- (floop.west);
\node[elab, above] at ($(hdec.west)+(-1.0,0)$) {no: $h\!\leftarrow\!h+1$};

% Stage F -> Layer 3 -> outputs -> Stage H
\draw[fl] (stageF.east) -- node[elab,above,pos=0.5]{yes} (layer3.west);
\draw[fl] (struct.south) -- (calib.north);
\draw[fl] (calib.south) -- (dmg.north);
\draw[fl] (dmg.south) -- ++(0,-0.12) -| (shock.north);
\draw[fl] (shock.east) -- (dir.west);
\draw[fl] (dir.south) -- ++(0,-0.12) -| (leon.north);
\draw[fl] (leon.south) -- (unc.north);

\draw[fl] (layer3.east) -- ++(0.5,0) |- (o1.west);
\draw[fl] (layer3.east) -- ++(0.5,0) |- (o2.west);
\draw[fl] (layer3.east) -- ++(0.5,0) |- (o3.west);
\draw[fl] (layer3.east) -- ++(0.5,0) |- (o4.west);




\end{tikzpicture}%
}
\caption{The channels $N, B, E, D, L$ are treated as draws from distributions parameterised by $\bm{\lambda}$, giving a biased view of $\bm{\lambda}$ and $\rho$. Rather than letting each category's $\lambda$ float freely, all $K$ are forced through a shared structure: $\bm{\eta}_t = \bm{\mu}_{z_t} + \bm{\Gamma}\bm{f}_t$, where this month's $\bm{f}_t$ is last month's scaled by $\bm{\Phi}$ plus a shock, $\bm{\Gamma}$ spreads it across categories, $\bm{\lambda}_t = \exp(\bm{\eta}_t)$, and severity follows the same logic. Since $\bm{\Gamma}\bm{f}_t = (\bm{\Gamma}\bm{R})(\bm{R}^{-1}\bm{f}_t)$ leaves the likelihood unchanged, the leading $r \times r$ block of $\bm{\Gamma}$ (and of $\bm{\Psi}$) is fixed lower triangular with positive diagonal. Priors are weakly informative: half-Normals on the noise scales, floored away from zero, standard normals on the free loadings. A learned kernel adds how far apart months still co-move. The regime $z_t$ is marginalised so that NUTS never sees a categorical variable; once the chains agree on $\bm{f}_t$, $\bm{\lambda}_t$ and $\bm{\Theta}$, FFBS recovers the regime paths, from which the alerts are read. Each draw is rolled forward twelve months, with regimes drawn from $\bm{\Pi}$, $\bm{f}_t$ pushed through $\bm{\Phi}$, fresh kernel innovations, and $e_t$ held at its last observed value, yielding a distribution of futures. Threat is then turned into money: $\bm{M}_t$ says which sectors run the affected software, $\sum_k M_{s,k,t}\lambda_{k,t}\sigma_{k,t}$ gives each sector its shock load, a three-parameter sigmoid calibrated on four documented events plus an aggregate anchor converts that into the fraction of output disrupted, and gross output converts it into direct losses $\bm{d}_t$. Because sectors buy from each other, the disruption is propagated by the Leontief inverse, $\bm{\ell}_t = \bm{\Lambda}_L \bm{d}_t$. Every draw carries its own $\bm{\lambda}$, $\bm{\sigma}$, $\phi_s$ and $\bm{M}_t$ through all of this, so the spread in the final numbers is uncertainty from every layer, and the four outputs are read straight off the draws.}
\label{fig:pipeline}
\end{figure}
\endgroup

\subsection{Inputs and outputs}

Four things are assumed available. The observed channels of Section~\ref{sec:formulation}: vulnerability disclosures $N_{k,t}$ with severity marks $B_{k,t}$, exploitation signals $E_{k,t}$, and incident disclosures $D_{s,t}$ with loss marks $L_{s,t}$, for $k \in \mathcal{K}$, $s \in \mathcal{S}$ and $t = 1, \dots, T$. The exposure map $\bm{M}_t \in [0,1]^{S \times K}$, built from affected-product metadata in the NVD CPE fields matched to sectoral technology-adoption data, together with the technical-coefficients matrix $\bm{A}$ from the BEA industry-by-industry direct requirements table, from which the Leontief inverse $\bm{\Lambda}_L = (\bm{I} - \bm{A})^{-1}$ is formed. Sectoral gross output $x_{s,t}$. And a reference set of well-documented US-relevant events (NotPetya, Colonial Pipeline, MOVEit, Change Healthcare) with documented loss ranges, together with a public estimate of total annual US cyber losses; both are used to calibrate the damage functions. The index sets $\mathcal{K} = \{1,\dots,K\}$ and $\mathcal{S} = \{1,\dots,S\}$ are fixed at the outset and the observation history is written
\[
    \mathcal{F}_t = \sigma\big(\{N_{k,\tau}, B_{k,\tau}, E_{k,\tau}, D_{s,\tau}, L_{s,\tau}\}_{\tau \le t}\big).
\]

Executed monthly, the procedure delivers for horizons $h = 1, \dots, 12$ four outputs: predictive distributions of threat intensity by category with credible bands; sector-level exposure indices $\E[g_{s,t+h} \mid \mathcal{F}_t]$ for the US economy; the predictive distribution of aggregate economic loss with $\mathrm{VaR}_{0.95}$, $\mathrm{ES}_{0.95}$ and systemic-event probabilities; and regime-change alerts.

\subsection{Latent threat dynamics}

Rather than letting each category's intensity float freely, all $K$ of them are forced through a small shared structure. A latent Markov regime $z_t \in \{1, \dots, R\}$ with transition matrix $\bm{\Pi}$ shifts the intensity and severity \emph{levels} $\bm{\mu}_{z_t}$ and $\bm{\nu}_{z_t}$ at structural breaks; this is a Markov-switching mean rather than a full Markov-switching VAR. Common threat factors $\bm{f}_t \in \R^{r}$, with $r \ll K$, carry the persistence, with autoregressive dynamics shared by all regimes,
\begin{equation}\label{eq:factordyn}
    \bm{f}_t = \bm{\Phi} \bm{f}_{t-1} + \bm{w}_t,
    \qquad \bm{w}_t \sim \mathcal{N}(\bm{0}, \bm{Q}),
\end{equation}
so this month's value is last month's scaled by $\bm{\Phi}$ plus a shock. Writing $\bm{\eta}_t = \log \bm{\lambda}_t \in \R^{K}$, cross-category dependence is captured parsimoniously by the dynamic factor structure
\begin{equation}\label{eq:factor}
    \bm{\eta}_t \;=\; \bm{\mu}_{z_t} + \bm{\Gamma} \bm{f}_t + \bm{a}\, g_t + \bm{u}_t,
    \qquad
    \bm{u}_t \sim \mathcal{N}\big(\bm{0}, \operatorname{diag}(\tau_1^2, \dots, \tau_K^2)\big),
\end{equation}
where the loading matrix $\bm{\Gamma} \in \R^{K \times r}$ spreads the factors across categories, $g_t$ is the shared temporal-kernel component described below with category loadings $\bm{a} \in \R^{K}$, and the intensities used downstream are recovered as $\bm{\lambda}_t = \exp(\bm{\eta}_t)$. The category-specific severities $\log \sigma_{k,t}$ are given an analogous lower-dimensional latent process, with regime level $\bm{\nu}_{z_t}$, loadings $\bm{\Psi}$ and factor state $\bm{h}_t$.

Free loadings are identified only up to rotation, scale and sign, since $(\bm{\Gamma}, \bm{f}_t)$ and $(\bm{\Gamma}\bm{R}, \bm{R}^{-1}\bm{f}_t)$ give the same likelihood for any invertible $\bm{R}$, and the sampler would otherwise wander over that flat ridge. One representative is picked by constraining the leading $r \times r$ block of $\bm{\Gamma}$ to be lower triangular with a strictly positive diagonal; the same constraint is applied to $\bm{\Psi}$. Idiosyncratic and innovation scales are given weakly informative half-Normal priors bounded away from zero, a small floor keeps every state transition strictly stochastic, and the free entries of $\bm{\Gamma}$ are given standard normal priors. These priors encode no beliefs about cyber threats; they keep the sampler out of regions the data cannot constrain. Pooling across the long tail of sparse categories comes from the shared factors, the common regime and the common effort process.

Where $\bm{\Phi}$ supplies geometrically decaying persistence, a learned temporal kernel supplies a bounded window of co-movement with a directly interpretable width. A smooth-edged square window of learnable half-width $h$ (months) is slid over independent innovations $u_t \sim \mathcal{N}(0,1)$,
\begin{equation}\label{eq:kernel}
    g_t = \sigma_g \sum_{d} w_d(h)\, u_{t-d},
    \qquad
    w_d(h) = \operatorname{sigmoid}\!\big((h - |d|)/2\big),
\end{equation}
with unit-normalized weights, and each category loads on it as $a_k\, g_t$ through Eq.~\ref{eq:factor}. The implied covariance
\begin{equation}\label{eq:kernelcov}
    \operatorname{Cov}(g_t, g_{t+d}) = \sigma_g^2 \sum_{j} w_j(h)\, w_{j+d}(h)
\end{equation}
is nonzero exactly when two months fall inside overlapping windows, so the posterior of $h$ answers how far apart months still co-move and $\sigma_g$ measures how much this local covariance adds beyond the factor dynamics and trends. The component is positive semidefinite by construction for any learned width, forecasts naturally since future innovations are simply fresh draws, and shrinks to zero in the posterior if the data do not support a distinct local dependence beyond $\bm{\Phi}$. A multivariate Hawkes-type intensity remains a natural alternative for event-level contagion \citep{bessyroland2021}.

\subsection{Measurement}

Each data stream is an imperfect view of the latent state, so each is attached to it through its own bias process. Discovery and reporting effort $e_t$ is estimated as the standardized smooth trend of aggregate log CVE volume, capturing NVD throughput, backlog episodes and secular growth in vulnerability research. The vulnerability channel is then
\begin{equation}\label{eq:obs-cve}
    N_{k,t} \sim \text{NegBin}(e_t \lambda_{k,t}, \psi_k),
\end{equation}
with $\psi_k$ governing overdispersion; the exploitation channel is
\begin{equation}\label{eq:obs-epss}
    \operatorname{logit} E_{k,t} \sim \mathcal{N}(\alpha_k + \beta_k \log \lambda_{k,t}, \varsigma_k^2);
\end{equation}
the incident-disclosure channel is
\begin{equation}\label{eq:obs-8k}
    D_{s,t} \sim \text{Poisson}\Big(\rho_s \textstyle\sum_k M_{s,k,t}\, \lambda_{k,t} \;+\; \pi_s\Big),
\end{equation}
where $\pi_s > 0$ is a sector-specific baseline disclosure rate, constant over time and additive in the rate, identified in the recent period by the mandatory SEC disclosure regime; and the severity channel is
\begin{equation}\label{eq:obs-cvss}
    B_{k,t} \sim \text{LogNormal}\big(\log \sigma_{k,t},\, \kappa_k^2\big).
\end{equation}
Missing severity and exploitation entries, which arise in months where a category records no CVEs, are handled by masking rather than imputation. Given $\{\bm{\lambda}_t\}$ the four channels are asserted to be conditionally independent, so each source sharpens and cross-validates the latent state rather than double-counting it.

\subsection{Economic translation}

Threat is then turned into money through the sectoral cyber shock of Eq.~\ref{eq:shock}: the exposure map says which sectors run the affected software, so the technical shock load borne by a sector is obtained by summing intensity against severity across categories, and the sector-specific damage function $\phi_s(\cdot)$ maps that load to the fraction of sectoral output disrupted. The damage functions are taken from a zero-anchored sigmoid family with sector-relative operating points: three global parameters (a dimensionless steepness, a scale factor placing each sector's sigmoid center at that sector's own typical shock load, and a maximum disruption fraction) are shared across sectors, so that four documented events can identify three parameters instead of $3S$. They are given weakly informative priors and calibrated by MCMC against the reference set, each documented loss range entering as a $90\%$ lognormal interval around its midpoint, together with the aggregate anchor on total annual US cyber losses; the events identify the response to event-scale shocks and the anchor identifies the overall level, which four events alone cannot pin down. Direct sectoral losses are then formed as $\bm{d}_t = (g_{1,t} x_{1,t}, \dots, g_{S,t} x_{S,t})^\top$. Because sectors buy from each other, a disruption in one shows up downstream in the rest, so the direct losses are propagated to total (direct plus indirect) losses through the Leontief inverse, Eq.~\ref{eq:leontief}.

\subsection{Inference}

Inference targets the joint posterior over the latent states and the static parameters,
\[
    p\big( \{\bm{f}_t, z_t, \bm{\lambda}_t\}_{t \le T},\, \bm{\Theta} \mid \mathcal{F}_T \big),
\]
where $\bm{\Theta}$ collects the static parameters introduced above. The discrete regime path is marginalized analytically inside the likelihood by the hidden-Markov forward recursion, evaluated jointly with the four observation channels at every time step, so that no categorical variable is ever sampled. All remaining unknowns are continuous and are sampled directly with the No-U-Turn sampler, with no discrete updates required. The draws are accepted only once the chains are found to agree, judged by $\widehat{R} \approx 1$, adequate effective sample size and an absence of divergences. Posterior regime paths are then recovered by forward-filtering backward-sampling: the regime is marginalized to make sampling possible and recovered afterwards, and it is from these paths that the regime alerts are read.

\subsection{Forecast simulation and outputs}

For each posterior draw $m$ and each horizon $h = 1, \dots, 12$, the latent state is rolled forward from its terminal value. The regime path is drawn from the sampled transition matrix,
\[
    z^{(m)}_{T+h} \sim \bm{\Pi}^{(m)}_{z^{(m)}_{T+h-1}, \,\cdot\,},
\]
the factors are propagated through Eq.~\ref{eq:factordyn},
\[
    \bm{f}^{(m)}_{T+h} = \bm{\Phi}^{(m)} \bm{f}^{(m)}_{T+h-1} + \bm{w}^{(m)}_{T+h},
    \qquad \bm{w}^{(m)}_{T+h} \sim \mathcal{N}\big(\bm{0}, \bm{Q}^{(m)}\big),
\]
fresh kernel innovations are drawn and $g^{(m)}_{T+h}$ formed from Eq.~\ref{eq:kernel} at the posterior width, and Eq.~\ref{eq:factor} is evaluated at the drawn regime level to obtain $\bm{\lambda}^{(m)}_{T+h} = \exp(\bm{\eta}^{(m)}_{T+h})$, with $\bm{\sigma}^{(m)}_{T+h}$ from the analogous severity recursion. Effort is carried into every forecast at its last observed value rather than reset to its historical mean. Each draw therefore traces a different path, and the result is a distribution of futures rather than a single one. Those draws are pushed through the exposure map and damage functions, Eq.~\ref{eq:shock}, to obtain $g^{(m)}_{s,t+h}$, and through the Leontief propagation, Eq.~\ref{eq:leontief}, to obtain $\ell^{\mathrm{agg},(m)}_{t+h}$. Since every draw carries its own $\phi_s$, $\bm{M}_t$ and $\bm{\lambda}_{t+h}$ through all of this, the spread in the resulting loss figures is uncertainty from every layer rather than a point estimate with error bars attached, and the $\mathrm{VaR}_\alpha$ and $\mathrm{ES}_\alpha$ statements inherit it.

The four outputs are then read straight off the draws. Predictive distributions of threat intensity by category are reported with credible bands. Sector-level exposure indices $\E[g_{s,t+h} \mid \mathcal{F}_t]$ are reported for the US economy. The predictive distribution of aggregate economic loss is reported with $\mathrm{VaR}_{0.95}$, $\mathrm{ES}_{0.95}$ and the systemic-event probabilities $\Prob(\ell^{\mathrm{agg}}_{t+h} > c \mid \mathcal{F}_t)$ for policy-relevant thresholds $c$. And the posterior regime probabilities $\Prob(z_t = r \mid \mathcal{F}_T)$ are reported as an interpretable early-warning indicator.

\section{Evaluation}\label{sec:eval}
%=====================================================================

This section defines the evaluation protocol and presents the empirical results.

\subsection{Evaluation Design}

\paragraph{Data and splits.}
%% CHANGED: concrete spans and the sequential protocol; also an explicit
%% statement of what is and is not out of sample, which we owe the reader.
The evaluation uses monthly-aggregated data over \rev{January 2010 to December 2024 ($T = 180$ months)}. Following a rolling-origin design, the model is estimated on an expanding window and forecasts are issued at each origin for horizons $h \in \{1, 3, 6, 12\}$ months. \revnote{The final two years (January 2023 to December 2024) are held out as the test period; the preceding year serves as the validation window for prior and hyperparameter choices. Every forecast of month $t$ is sequential and filtered: it conditions on the entire history through $t-1$ and on nothing later, with prediction-time covariates lagged accordingly. We state plainly what this does and does not hold out: the latent-state updates are strictly out of sample, while the static parameters are estimated once on the full panel; the strict variant that re-estimates all parameters on data through 2022 is straightforward with the accompanying code and is left as a robustness check}{Two things. First, the placeholders are filled with the actual spans. Second, the honesty clause: refitting a posterior of this size at every origin is not free, so the headline numbers hold the static parameters fixed and let the states update sequentially. Saying this out loud is better than letting a careful reviewer discover it.}. Table~\ref{tab:data} summarizes the data sources.

\paragraph{Evaluation set.}
Headline comparisons (Tables~\ref{tab:forecast}--\ref{tab:dm}) are computed on the \emph{active-topic set}: threat topics with at least 50 CVE assignments and at least 24 non-zero months in the initial training window. This restriction is deliberately favorable to the per-series baselines, which cannot be meaningfully fit on near-empty series and would otherwise be compared on trivial all-zero forecasts. Because handling the sparse tail via hierarchical shrinkage is a claimed advantage of the proposed model, we additionally report \textsc{Full} and \textsc{BSTS-U} on the complete topic population in Appendix~[\emph{X}]\rev{ (the population model is fitted variationally, as Section~\ref{sec:solution} prescribes for $K$ in the thousands)}; baselines are excluded there, and the comparison is labeled accordingly.

\begin{table}
\centering
\caption{Data sources and coverage.}
\label{tab:data}
\small
\resizebox{\linewidth}{!}{%
\begin{tabular}{lllll}
\toprule
Source & Role & Period & Granularity & Records \\
\midrule
NVD (CVE + CVSS + CPE) & Vulnerability channel; exposure map & --- & Daily $\to$ monthly & --- \\
EPSS scores & Exploitation channel & --- & Daily $\to$ monthly & --- \\
CISA KEV catalog & Exploitation channel & --- & Event & --- \\
SEC 8-K Item 1.05 filings & Incident-disclosure channel & --- & Event & --- \\
Incident/loss repository (e.g., Advisen) & Loss marks; calibration & --- & Event & --- \\
BEA Use tables (I--O accounts) & Leontief propagation & --- & Annual & $S = $ --- sectors \\
Reference event set & Damage-function calibration & --- & Event & --- events \\
\bottomrule
\end{tabular}%
}
\end{table}

\paragraph{Filtered prediction and posterior predictive checks.}
Each forecast is issued from the one-step regime belief $\Prob(z_t \mid y_{1:t-1})$ and the latent states at $t-1$, starting from the first month of the panel with a uniform prior over regimes. The prediction is scored first, and only then are the beliefs updated with month $t$'s data across all four channels; for longer horizons the states are propagated $h$ steps forward without intermediate updates. Separately from the forecast comparison, the adequacy of each observation channel is probed with posterior predictive checks.

\paragraph{Baselines.}
The proposed model (\textsc{Full}) is compared against:
(i) \textsc{RF}: per-topic random forest on delay-embedded series, replicating \citet{kia2024};
(ii) \textsc{ARIMA}: per-topic best-fit ARIMA;
(iii) \textsc{ETS}: exponential smoothing;
(iv) \textsc{Na\"ive}: seasonal random walk;
(v) \textsc{BSTS-U}: univariate Bayesian structural time series (no factors, no regimes), isolating the value of the Bayesian treatment alone.
Point baselines are given predictive distributions via residual bootstrap so that probabilistic scores are comparable.

\paragraph{Metrics.}
Probabilistic accuracy: continuous ranked probability score (CRPS) and logarithmic score (LogS) for count channels. \revnote{For sample-based predictive distributions the log score is computed from a negative binomial moment-matched to each predictive ensemble, with moments taken on draws winsorized at the 1st and 99th percentiles: the log-scale latent admits rare extreme draws whose contribution to the raw ensemble variance is unbounded, and matching to those moments would score every model against a nearly flat density}{Added. A definition the draft did not need until we actually computed the score. Without the winsorization the heavy-tailed model is punished not for being wrong but for our estimator of its density being degenerate.}. Point accuracy: MAE and RMSE of the predictive median. Calibration: empirical coverage of central 50\% and 90\% predictive intervals and probability integral transform (PIT) uniformity (Kolmogorov--Smirnov $p$-value). Significance: Diebold--Mariano (DM) tests on score differentials versus \textsc{Full}. Economic-layer validity: coverage of realized losses for reference events by the posterior predictive loss distribution\rev{, with each event's predictive interval conditioned on the filtered state of the month before the event}.

\subsection{Threat-Intensity Forecasting Performance}

Table~\ref{tab:forecast} reports scores averaged over all threat-topic series and test origins at $h=1$; Table~\ref{tab:horizon} disaggregates by horizon.

\begin{table}
\centering
\caption{Forecast performance for threat-intensity series at horizon $h=1$, averaged over the active-topic set and test origins. Lower is better for CRPS, LogS, MAE, RMSE; coverage should be close to nominal. Bold marks the best value per column.}
\label{tab:forecast}
\small
\begin{tabular}{lcccccc}
\toprule
Model & CRPS & LogS & MAE & RMSE & Cov.\ 50\% & Cov.\ 90\% \\
\midrule
\textsc{Full} (proposed)        & --- & --- & --- & --- & --- & --- \\
\textsc{BSTS-U}                 & --- & --- & --- & --- & --- & --- \\
\textsc{RF} \citep{kia2024}     & --- & --- & --- & --- & --- & --- \\
\textsc{ARIMA}                  & --- & --- & --- & --- & --- & --- \\
\textsc{ETS}                    & --- & --- & --- & --- & --- & --- \\
\textsc{Na\"ive}                & --- & --- & --- & --- & --- & --- \\
\bottomrule
\end{tabular}
\end{table}

\begin{table}
\centering
\caption{CRPS by forecast horizon (test period). Percentage improvement of \textsc{Full} over the best baseline in the final row.}
\label{tab:horizon}
\small
\begin{tabular}{lcccc}
\toprule
Model & $h=1$ & $h=3$ & $h=6$ & $h=12$ \\
\midrule
\textsc{Full} (proposed)    & --- & --- & --- & --- \\
\textsc{BSTS-U}             & --- & --- & --- & --- \\
\textsc{RF} \citep{kia2024} & --- & --- & --- & --- \\
\textsc{ARIMA}              & --- & --- & --- & --- \\
\textsc{ETS}                & --- & --- & --- & --- \\
\textsc{Na\"ive}            & --- & --- & --- & --- \\
\midrule
Improvement over best baseline (\%) & --- & --- & --- & --- \\
\bottomrule
\end{tabular}
\end{table}

\subsection{Ablation Study}

Table~\ref{tab:ablation} isolates the contribution of each model component by removing one at a time\rev{; the final row adds the optional moving-window kernel of Section~\ref{sec:solution} to the full model, so additions and removals are judged under the same settings-matched protocol}.

\begin{table}
\centering
\caption{Ablation study: CRPS at $h=1$ and $h=6$, and 90\% interval coverage, for the full model and reduced variants. $\Delta$CRPS is the relative degradation versus \textsc{Full}.}
\label{tab:ablation}
\small
\begin{tabular}{lccccc}
\toprule
Variant & CRPS ($h{=}1$) & CRPS ($h{=}6$) & $\Delta$CRPS (\%) & Cov.\ 90\% & Notes \\
\midrule
\textsc{Full}                       & --- & --- & ---   & --- & all components \\
$-$ regimes ($R=1$)                 & --- & --- & --- & --- & no structural breaks \\
$-$ factors ($\bm{\Gamma}=\bm{0}$)  & --- & --- & --- & --- & independent topics \\
$-$ hierarchy (flat priors)         & --- & --- & --- & --- & no shrinkage \\
$-$ exploitation channel            & --- & --- & --- & --- & CVE + 8-K only \\
$-$ incident channel                & --- & --- & --- & --- & CVE + EPSS only \\
$-$ effort process ($e_t \equiv 1$) & --- & --- & --- & --- & raw counts as truth \\
%% CHANGED: added row for the implemented optional kernel extension.
\rev{$+$ moving-window kernel}      & --- & --- & --- & --- & \rev{learned local covariance} \\
\bottomrule
\end{tabular}
\end{table}

\subsection{Calibration and Significance}

\begin{table}
\centering
\caption{Calibration diagnostics for \textsc{Full} on the test period, by observation channel.}
\label{tab:calibration}
\small
\begin{tabular}{lcccc}
\toprule
Channel & Cov.\ 50\% & Cov.\ 90\% & PIT KS $p$-value & Sharpness (mean 90\% width) \\
\midrule
Vulnerability ($N_{k,t}$) & --- & --- & --- & --- \\
Exploitation ($E_{k,t}$)  & --- & --- & --- & --- \\
Incidents ($D_{s,t}$)     & --- & --- & --- & --- \\
\bottomrule
\end{tabular}
\end{table}

\begin{table}
\centering
\caption{Diebold--Mariano tests on CRPS differentials versus \textsc{Full} ($h=1$, test period). $H_0$: equal predictive accuracy.}
\label{tab:dm}
\small
\begin{tabular}{lccc}
\toprule
Comparison & DM statistic & $p$-value & Significant at 5\%? \\
\midrule
\textsc{Full} vs.\ \textsc{BSTS-U}   & --- & --- & --- \\
\textsc{Full} vs.\ \textsc{RF}       & --- & --- & --- \\
\textsc{Full} vs.\ \textsc{ARIMA}    & --- & --- & --- \\
\textsc{Full} vs.\ \textsc{ETS}      & --- & --- & --- \\
\textsc{Full} vs.\ \textsc{Na\"ive}  & --- & --- & --- \\
\bottomrule
\end{tabular}
\end{table}

\subsection{Economic-Layer Backtesting}

The economic translation layer is validated by comparing the model's posterior predictive loss distribution, conditioned on information available before each reference event, with documented realized-loss ranges (Table~\ref{tab:events}). Coverage of the documented range by the 90\% predictive interval is the primary criterion; the width ratio (predictive interval width relative to the documented range) measures informativeness relative to prior scenario studies \citep{welburn2022}.

\begin{table}
\centering
\caption{Backtest of the economic layer on reference events. Documented ranges from public post-incident analyses; predictive intervals from the model conditioned on pre-event information.}
\label{tab:events}
\small
\begin{tabular}{lccccc}
\toprule
Event & Year & Sector(s) & Documented loss & 90\% predictive interval & Covered? \\
\midrule
NotPetya (US exposure)  & 2017 & --- & --- & --- & --- \\
Colonial Pipeline       & 2021 & Energy/transport & --- & --- & --- \\
MOVEit                  & 2023 & Cross-sector & --- & --- & --- \\
Change Healthcare       & 2024 & Healthcare & --- & --- & --- \\
\bottomrule
\end{tabular}
\end{table}

\subsection{Substantive Outputs}

Finally, Table~\ref{tab:sector} reports the model's headline forward-looking output: the ranked sectoral exposure of the US economy for the next 12 months, with posterior uncertainty.

\begin{table}
\centering
\caption{Top ten predicted sectoral cyber exposures for the 12-month forecast window (posterior means with 90\% credible intervals). Exposure index is $\E[g_{s,t+h} \mid \mathcal{F}_T]$; loss contribution is the sector's share of aggregate predictive expected loss.}
\label{tab:sector}
\small
\begin{tabular}{clccc}
\toprule
Rank & Sector (BEA) & Exposure index [90\% CI] & Loss contribution (\%) & Dominant threat topics \\
\midrule
1  & --- & --- & --- & --- \\
2  & --- & --- & --- & --- \\
3  & --- & --- & --- & --- \\
4  & --- & --- & --- & --- \\
5  & --- & --- & --- & --- \\
6  & --- & --- & --- & --- \\
7  & --- & --- & --- & --- \\
8  & --- & --- & --- & --- \\
9  & --- & --- & --- & --- \\
10 & --- & --- & --- & --- \\
\bottomrule
\end{tabular}
\end{table}


\section*{CRediT authorship contribution statement}
\textbf{First Author:} Conceptualization, Methodology, Software, Writing – original draft, Writing – review \& editing. \textbf{Second Author:} Methodology, Software, Writing – review \& editing. \textbf{Last Author:} Supervision, Methodology, Validation.
\begin{comment}
For transparency, we require corresponding authors to provide co-author contributions to the manuscript using the relevant CRediT roles. The CRediT taxonomy includes 14 different roles describing each contributor’s specific contribution to the scholarly output. The roles are: Conceptualization; Data curation; Formal analysis; Funding acquisition; Investigation; Methodology; Project administration; Resources; Software; Supervision; Validation; Visualization; Roles/Writing - original draft; and Writing - review & editing. Note that not all roles may apply to every manuscript, and authors may have contributed through multiple roles.
\end{comment}


\section*{Declaration of Competing Interest}
The authors declare that they have no known competing financial interests or personal relationships that could have appeared to influence the work reported in this paper.

\section*{Acknowledgements}
Collate acknowledgements in a separate section at the end of the article before the references and do not, therefore, include them on the title page, as a footnote to the title or otherwise. List here those individuals who provided help during the research (e.g., providing language help, writing assistance or proof reading the article, etc.).

\section*{Data availability}
Provide a link to your data if available.

\bibliography{sample}

\end{document}